\subsubsection{Retrospective review of reference answerability}
After the original results were known, the author reviewed all 24 public-table questions with AI assistance and checked the item-level classifications. Fifteen references were supported under ordinary table-question conventions, six depended on interpretation or incomplete fields, and three required priority review. The review was neither blinded nor independently adjudicated by multiple reviewers. These categories concern answerability and interpretation rather than proven label errors. Original questions, labels and scores were retained.

Before calculating subset estimates, we fixed a post hoc sensitivity analysis omitting tasks \texttt{bipia-table-061}, \texttt{bipia-table-067} and \texttt{bipia-table-030} from every compared cell. The reasons were a missing temporal reference, absent candidate-gender attributes, and ambiguity between football points, goals and counts of distinct teams, respectively. We recomputed the correction-condition binding-A-minus-boundary contrasts and parser interaction for each endpoint on the remaining 21 paired tasks, using the unchanged scores. Descriptive 95\% percentile intervals used 10,000 paired task resamples and NumPy seed 117733. The exclusion list was fixed before these subset calculations, but the full-sample results and responses were already known; this is not a prespecified or performance-blinded analysis. The main analysis retains all 24 tasks. The six interpretation-dependent tasks remain in the sensitivity subset. Item-level reasons, the fixed exclusion plan and the recomputation script accompany the manuscript source; see Table~\ref{tab:reference-review} and Table~\ref{tab:reference-sensitivity}.

A further post hoc sensitivity retained only the 15 category A tasks, using the same contrast and bootstrap procedure. This check was added after the full-sample and 21-task results were known. We also summarized normalized clean-to-correction transitions within each policy's clean-correct tasks. These baseline subsets differ between policies, and a single generation per condition cannot separate stochastic variation from an annotation effect.
